%% Style by Bram Gotink, as a clone of the LyX style orginally used, but in clean code.
\documentclass[dutch]{article}
\usepackage[T1]{fontenc}
\usepackage[utf8]{inputenc}
% XeLaTeX/Tectonic negeren inputenc: zet deze tekens om, anders verdwijnen ze
% of worden ze fout gedrukt (§ werd ğ). Onder pdflatex verandert er niets.
\usepackage{iftex}
\ifTUTeX
  \usepackage{newunicodechar}
  \newunicodechar{§}{\S}
  \newunicodechar{‘}{\textquoteleft}
  \newunicodechar{’}{\textquoteright}
  \newunicodechar{“}{\textquotedblleft}
  \newunicodechar{”}{\textquotedblright}
\fi
\usepackage{geometry}
\geometry{verbose,letterpaper,tmargin=2.5cm,bmargin=2.5cm,lmargin=2.5cm,rmargin=2.5cm}
\setlength{\parskip}{\medskipamount}
\setlength{\parindent}{0pt}
\usepackage{textcomp}
\usepackage{titlesec}
\usepackage{enumitem}
\usepackage{color}
\usepackage{ulem}
\usepackage{hyperref}
%\hypersetup{colorlinks=false,urlcolor=cyan,pdfborder=0 0 0}

% VTK Styling

% subsections worden gebruikt als artikels in het HR en Verkiezingsreglement,
% Deze nummering moet doorlopen
\counterwithout{subsection}{section}

% parts zijn "Titel I - De vereniging"
\def\thepart{\Roman{part}}
\def\labelpart{Titel}
\titleformat{\part}
            [hang]
            {\normalfont\huge\bfseries}
            {\labelpart~\thepart~-}
            {10pt}
            {}
% sections zijn "Artikel 1. Naam" enzo
\def\thesection{\arabic{section}}
\def\labelsection{Artikel}
\titleformat{\section}
            [hang]
            {\normalfont\Large\bfseries}
            {\labelsection~\thesection.}
            {10pt}
            {}
% subsections zijn "Artikel 1. Naam", enkel gebruikt in HR
% (daar zijn sections "Afdeling 1 - Naam", zie de .tex van HR)
\def\thesubsection{\arabic{subsection}}
\def\labelsubsection{}
\titleformat{\subsection}
			[hang]
			{\normalfont\large\bfseries}
			{\labelsubsection~\thesubsection.}
			{10pt}
			{}
% hoogste niveau lijsten zijn §1. etc
\def\theenumi{\S\arabic{enumi}}
\def\labelenumi{\theenumi.}
% tweede niveau wordt eigenlijke eerste niveau
\def\theenumii{\arabic{enumii}}
\def\labelenumii{\theenumii.}
% eerste niveau itemize komt enkel voor in een lijst
\def\labelitemi{}
% layout de lijsten wat mooier
\setlist[enumerate,1]{align=left,labelindent=0pt,labelsep=3em,leftmargin=5em}
\setlist[enumerate,2]{align=left}
\setlist[itemize,1]{align=left,labelindent=0pt,labelsep=3em,leftmargin=5em,labelwidth=0pt}

% locate changes and color them red.
\newenvironment{aanpassing}{\color{red}}{}
\newcommand{\aangepast}[1]{\begin{aanpassing}#1\end{aanpassing}}
\newcommand{\verwijderd}[1]{\aangepast{\sout{#1}}}
% nieuwe (voorgestelde) tekst in het groen
\definecolor{vtkgroen}{rgb}{0,0.55,0}
\newcommand{\toegevoegd}[1]{{\color{vtkgroen}#1}}
\newcommand{\draft}{\textit{\textbf{draft}}}

\usepackage{babel}



\begin{document}


\title{Huishoudelijk Reglement Vlaamse Technische Kring vzw}
\date{30 juni 2025}

\maketitle

\part{Algemene bepalingen}\label{art:alg_bepalingen}
\section{Algemene Bepalingen}
\begin{enumerate}
    \item Dit Huishoudelijk Reglement is geschreven naar aanleiding van de statuten van de Vlaamse Technische Kring vzw.
    \item Indien een lid van de Algemene Vergadering dit reglement niet naleeft zonder toestemming van het Bestuursorgaan, kan deze persoon persoonlijk aansprakelijk gesteld worden voor de gevolgen van zijn overtreding.
    \item Verder in dit document wordt de handelsnaam VTK gebruikt voor zowel Vlaamse Technische Kring vzw, als de faculteitskring VTK als het facultair overlegorgaan VTK.
    \item De Algemene Vergadering kan met gewone meerderheid van de aanwezige of vertegenwoordigde leden beslissen over wijzigingen in dit document. Als het gaat over een aanpassing van Titel 3 in het Huishodelijk Reglement, moeten aanpassingen eerst besproken worden op een vergadering waar alle voorzitters uitgenodigd zijn alvorens deze op een AV wordt gebracht.
\end{enumerate}

\part{Organen en functies}
\section{Bestuurders}
\begin{enumerate}
    \item Het Bestuursorgaan kiest onder de bestuurders een voorzitter, rekening houdend met de voorwaarden in \ref{BO_voorwaardenVoorzitter} van dit artikel, ondervoorzitter, secretaris en minstens één penningmeester. 
    \item\label{BO_voorwaardenVoorzitter} Het Bestuursorgaan zal bij het kiezen van een voorzitter rekening houden met volgende voorwaarde: 
        \begin{itemize}
        \item De voorzitter moet in het jaar dat zijn mandaat van kracht is, zijn bachelordiploma aan de Faculteit Ingenieurswetenschappen behaald hebben en ingeschreven zijn voor een masteropleiding aan de Faculteit Ingenieurswetenschappen.
        \end{itemize}
    \item De bestuurders voeren hun mandaat onbezoldigd uit.
\end{enumerate}

\section{Bijeenkomst van het Bestuursorgaan}
    \begin{enumerate}
    \item Het Bestuursorgaan vergadert na oproeping door de voorzitter zo dikwijls als het belang van de vereniging het vereist. De vergadering wordt gehouden op de plaats die in het oproepingsbericht wordt vermeld.
    \item De oproeping dient minstens 24 uur voor aanvang van de vergadering te gebeuren.
    \item Het Bestuursorgaan wordt voorgezeten door de voorzitter. Indien deze niet aanwezig kan zijn stelt hij een vervanger aan. Deze vervanger wordt vermeld in de uitnodiging. In het geval er geen vervanger aangesteld is wordt het Bestuursorgaan voorgezeten door de ondervoorzitter. Indien de ondervoorzitter eveneens afwezig is, wordt het Bestuursorgaan voorgezeten door de oudst aanwezige dagelijks bestuurder. Indien deze allen afwezig zijn, beslist het Bestuursorgaan onder leiding van de oudst aanwezige bestuurder of de bijeenkomst wordt uitgesteld.
    \item Het Bestuursorgaan kan slechts geldig vergaderen wanneer ten minste 1/4 van zijn leden aanwezig zijn. Het Bestuursorgaan oefent zijn bevoegdheden uit als college.
    \item Van elke vergadering worden notulen opgesteld, ondertekend door de voorzitter en de secretaris van het Bestuursorgaan. De notulen kunnen tevens getekend worden door de bestuurders die het wensen. Ze worden bewaard in een notulenregister dat ter inzage zal zijn van de werkende leden, die hun inzagerecht zullen uitoefenen overeenkomstig de in de WVV-wet vastgelegde modaliteiten.
    \item Bij iedere bijeenkomst van het Bestuursorgaan dient de Kringvergadering uitgenodigd worden.
    \end{enumerate}

\section{Belangenvermenging}
    \begin{enumerate}
        \item Indien een bestuurder, rechtstreeks of onrechtstreeks, voor een beslissing die onder de bevoegdheid valt van het Bestuursorgaan een belang van vermogensrechtelijke of persoonlijke aard heeft dat strijdig is met het belang van de vereniging, moet hij dit meedelen aan de andere bestuurders voor het Bestuursorgaan een besluit neemt. De verklaring van deze bestuurder over de aard van het strijdig belang wordt opgenomen in de notulen van de vergadering.
        \item Het Bestuursorgaan kan met gewone meerderheid van de stemmen beslissen of iemand al dan niet een belangenconflict van vermogensrechtelijke of persoonlijke aard heeft.
        \item De bestuurder met het tegenstrijdig belang verwijdert zich uit de vergadering en onthoudt zich van de beraadslaging en de stemming over de aangelegenheid waarop het betrekking heeft.
        \item Voormelde procedure is niet toepasselijk op gebruikelijke verrichtingen die plaatshebben onder de voorwaarden en tegen de zekerheden die op de markt gewoonlijk gelden voor soortgelijke verrichtingen.
        \item Indien een bestuurder lid is van een (niet-)verkiesbare ploeg verwijdert hij zich uit de vergadering indien er een beraadslaging plaatsvindt die betrekking heeft op de controle op de verkiezingen of de bespreking van het verkiezingsresultaat.
    \end{enumerate}

 \section{Bevoegdheden van het Bestuursorgaan}
    \begin{enumerate}
    \item Het Bestuursorgaan is bevoegd om alle handelingen van intern bestuur te verrichten die nodig of dienstig zijn tot verwezenlijking van het doel van de vereniging, met uitzondering van die handelingen waarvoor volgens de wet, de statuten en/of dit Huishoudelijk Reglement de Algemene Vergadering exclusief bevoegd is.
    \item Het Bestuursorgaan beslist over het Intern Reglement bij gewone meerderheid van de stemmen van de aanwezige of vertegenwoordigde werkende leden. Voorafgaande aan de beslissing is een advies van de Kringvergadering vereist.
    \item Onverminderd de verplichtingen die voortvloeien uit collegiaal bestuur, met name overleg en toezicht, kunnen de bestuurders de bestuurstaken onder elkaar verdelen. Zodanige verdeling van taken kan aan derden niet worden tegengeworpen, zelfs niet nadat ze zijn openbaar gemaakt. Niet-naleving ervan brengt wel de interne aansprakelijkheid van de betrokken bestuurder(s) in het gedrang.
    \item Het Bestuursorgaan kan een deel van zijn bestuursbevoegdheid delegeren aan één of meerdere derden niet-bestuurders, zonder dat deze overdracht evenwel betrekking kan hebben op het algemeen beleid van de vereniging of de algemene bestuursbevoegdheid van het Bestuursorgaan. Hierin wordt een onderscheid gemaakt tussen vertegenwoordigers van werkgroepen en bijzondere mandaten. De bevoegdheid van vertegenwoordigers van werkgroepen wordt beschreven in dit Huishoudelijk Reglement, die van andere mandaten wordt precies afgebakend door het Bestuursorgaan, die ook de duur van het mandaat bepaalt. Het mandaat kan te allen tijde met onmiddellijke ingang door het Bestuursorgaan worden ingetrokken.
    \end{enumerate}

\section{De plichten van het Bestuursorgaan}
    \begin{enumerate}
    \item De voorzitter van het Bestuursorgaan nodigt de leden van de Kringvergadering, bestaande uit alle verkozen praesidiumleden in een bepaald werkingsjaar, uit als waarnemer.
    \item De voorzitter van het Bestuursorgaan stuurt de verslagen van het Bestuursorgaan door naar de leden van de Kringvergadering.
    \item Het Bestuursorgaan bekrachtigt alle beslissingen van de Kringvergadering. Indien het Bestuursorgaan een beslissing van de Kringvergadering niet bekrachtigt, dient de voorzitter van het Bestuursorgaan alle leden van de Kringvergadering schriftelijk op de hoogte te stellen van de genomen beslissing en deze beslissing te motiveren.
    % \item Het Bestuursorgaan is verplicht elk jaar tussen februari en juni verkiezingen te organiseren, de "praesidiumverkiezingen"\ genaamd. Op deze verkiezingen zijn alle studenten aan de Faculteit Ingenieurswetenschappen aan de KU Leuven stemgerechtigd. De precieze modaliteiten van deze verkiezingen moeten worden vastgelegd in een specifiek verkiezingsreglement, dat door het Bestuursorgaan moet worden goedgekeurd. De verkiezingen loopt eveneens met het oog op de verkiezingen van de vertegenwoordigers van de studenten bij de KU Leuven. Daarbij respecteert het Bestuursorgaan het Participatiedecreet en het Participatiereglement van Stura KU Leuven en LOKO.
    % \item Het Bestuursorgaan ziet er op toe dat de verkiezingen correct verlopen. Vervolgens moet het Bestuursorgaan beslissen of de verkiezingen correct zijn verlopen, dit Huishoudelijk Reglement en het verkiezingsreglement indachtig. De bekrachtiging van het verkiezingsresultaat houdt in dat de leden van de winnende ploeg op de laatste Algemene Vergadering van het werkingsjaar toetreden als werkend lid.
    % \item Indien het verkiezingsresultaat niet bekrachtigd wordt, dienen de dagelijks bestuurders die het komende academiejaar niet langer aan de Faculteit Ingenieurswetenschappen zullen zijn een aantal personen aan te wijzen die als werkend lid tot de vereniging worden toegelaten. Deze personen staan in voor de continuïteit van de werking van VTK en zijn verplicht nieuwe verkiezingen te organiseren. 
    \end{enumerate}

\section{Vertegenwoordiging}\label{ArtikelVertegenwoordiging}
    \begin{enumerate}
    \item Het Bestuursorgaan vertegenwoordigt als college de vereniging in alle handelingen in en buiten rechte.
    \item Onverminderd de algemene vertegenwoordigheidsbevoegdheid van het Bestuursorgaan als college, wordt de vereniging in en buiten rechte eveneens geldig vertegenwoordigd door ofwel:
        \begin{itemize}
        \item twee dagelijks bestuurders, waarvan één van beiden voorzitter, ondervoorzitter, penningmeester of secretaris is; ofwel
        \item één dagelijks bestuurder tot een bedrag van EUR 2500,- voor wat het dagelijks bestuur betreft dat binnen de bevoegdheid valt van deze persoon zoals beschreven in het Intern Reglement.
        \end{itemize}\label{ParLimietenVertegenwoordiging}
    \item Het Bestuursorgaan kan, wanneer de noodzakelijkheid en het belang van de vereniging zulks vereisen, personen uitzonderlijk machtiging geven om ten aanzien van derden te handelen tot een door het Bestuursorgaan gespecificeerd bedrag.
    \end{enumerate}

\section{De Grocovergadering}
\begin{enumerate}
    \item In de Grocovergadering, ook wel Grocomeet genoemd, zetelen de dagelijks bestuurders. Dit omvat de groepsco\"ordinatoren en Groep 5.
    \item De Grocovergadering bereidt de Kringvergadering voor. De Praeses roept de Kringvergadering samen, maar de Grocovergadering stelt de agenda op en is verantwoordelijk voor de uitvoering van de besluiten van de Kringvergadering.
    \item Het dagelijks bestuur wordt waargenomen door de Grocovergadering.
    \item De Grocovergadering heeft slechts een adviserende rol aan de Kringvergadering met betrekking tot budgettaire beslissingen. 
    \end{enumerate}

\section{De Kringvergadering}
\begin{enumerate}
    \item De Kringvergadering bestaat uit alle leden die verkozen zijn volgens het participatiereglement van LOKO, hierna het Praesidium genoemd.
    \item Dit Huishoudelijk Reglement, in samen met het Intern Reglement, vormt eveneens het kringreglement van de faculteitskring VTK, waar nodig aangevuld met het verkiezingsreglement. 
    \item De Kringvergadering geeft advies aan het Bestuursorgaan over financi\"ele beslissingen, alsook over wijzigingen van het Intern Reglement en het verkiezingsreglement.
    \item Een lid van de Kringvergadering dat ontslagen is als werkend lid, is automatisch ook ontslagen als lid van de Kringvergadering. 
    \item De beslissingen van de Kringvergadering worden genomen via consensus, indien dit niet mogelijk
is wordt er gestemd. Alle leden van de Kringvergadering hebben gelijk stemrecht. Elk lid heeft
één stem.
    \item De leden van de Kringvergdering kunnen te allen tijde afgezet worden door de Algemene Verga-
dering die daarover beslist bij 2/3 meerderheid van de stemmen van de aanwezige of vertegen-
woordigde werkende leden
    \item Leden kunnen ook zelf ontslag nemen door een ontslagbrief te richten aan de Kringvergadering of
aan de voorzitter van de Kringvergadering.
    \item Een lid die door de omstandigheden geen werkend lid van de vereniging meer kan zijn, wordt van
rechtswegen afgezet als lid van de Kringvergadering.
\end{enumerate}

\section{Het facultair overlegorgaan}
\begin{enumerate}
    \item De Kringvergadering van VTK geldt als de Algemene Vergadering van het facultair overlegorgaan VTK, hierna het FO genoemd. 
    \item Praesidiumleden aangesteld door de Algemene Vergadering die niet verkozen zijn conform de bepalingen van het verkiezingsreglement nemen een adviserende rol waar.
    \item Het FO is bevoegd voor het verkiezen van studentenvertegenwoordigers in de opleiding en departementale organen en voor de vertegenwoordiging van de student op facultair (FPOC, faculteitsraad en faculteitsbestuur) en universitair (Stura AV) niveau.
    \item Dit Huishoudelijk Reglement vormt eveneens het Intern Reglement van het FO VTK, waar nodig aangevuld met het Intern Reglement en het verkiezingsreglement.
\end{enumerate}

\section{Het Auditcomit\'e}
    \subsection{Samenstelling}
    \begin{enumerate}
    \item Het Auditcomit\'e bestaat uit minstens drie personen, zodat representatie op de Algemene Vergadering verzekerd wordt. Het lidmaatschap van het Auditcomit\'e is onverenigbaar met het lidmaatschap van het Bestuursorgaan en/of het Dagelijks Bestuur.
    \item De Praeses draagt elk werkingsjaar personen voor voor het Auditcomit\'e. Deze worden door de Algemene Vergadering verkozen voor een termijn van \'e\'en jaar. Zij kunnen onbeperkt herkozen worden.
    \item De Algemene Vergadering kan steeds een lid van het Auditcomit\'e ontslaan, volgens dezelfde regels als voor de uitsluiting van een Werkend Lid, of een lid van het Auditcomit\'e kan steeds zijn ontslag indienen.
    \item Het Auditcomit\'e kiest uit hun midden een voorzitter voor het komende jaar. Dit mag niet afdoen aan zijn collegiaal karakter. De voorzitter van het Auditcomit\'e is verantwoordelijk voor het samenroepen van vergaderingen en rapportering aan het Bestuursorgaan. 
    \item Het Auditcomit\'e kan leden van het Bestuursorgaan uitnodigen voor zijn vergaderingen.
    \end{enumerate}

    \subsection{Doel}
    \begin{enumerate}
    \item Het Auditcomit\'e heeft tot doel een externe controle te bieden aan het Bestuursorgaan op het gebied van financi\"ele en administratieve verplichtingen. Hierbij handelen ze objectief en zonder de eigenheid van het bestuur te schaden.
    \item Het Auditcomit\'e is een adviesorgaan. Zijn belangrijkste taak is het informeren van het Bestuursorgaan over zijn verplichtingen en te rapporteren aan de Algemene Vergadering over de uitvoering van deze verplichtingen.
    \end{enumerate}

    \subsection{Opdracht}
    \begin{enumerate}
    \item Het Auditcomit\'e zal een advies geven over de technische correctheid en haalbaarheid van de begroting. Het Bestuursorgaan bezorgt hiertoe tijdig een exemplaar van de begroting.
    \item Na de controle op het einde van het eerste semester geeft het Auditcomit\'e opnieuw een advies over de begroting. Hiertoe krijgen zij tijdig de beschikbare resultaten van eerste semester.
    \item Het Auditcomit\'e zal doorheen het jaar opvolgen of het Bestuursorgaan de boekhouding correct voert en voldaan wordt aan alle verplichtingen (zoals btw, rechtspersonenbelasting, publicaties staatsblad).
    \item Het Auditcomit\'e zal na afloop van het boekjaar de boekhouding controleren en een advies geven over het al dan niet goedkeuren van de jaarrekening en het al dan niet kwijten van het Bestuursorgaan.
    \item Het Auditcomit\'e kan steeds in bijzondere gevallen een advies formuleren aan de Algemene Vergadering.
    \item Elk orgaan van VTK kan te allen tijde een advies vragen aan het Auditcomit\'e.
    \item De adviezen vermeld in \S1 en \S4 van dit artikel zijn vereist vooraleer de Algemene Vergadering tot een beslissing kan overgaan, zoals bepaald in de statuten van VTK. 
    \item De adviezen/controles vermeld in \S2 en \S4 moeten verplicht aangevraagd/uitgevoerd worden, maar het ontbreken ervan heeft geen gevolgen voor de beslissingsmogelijkheden op de Algemene Vergadering.
    \end{enumerate}

    \subsection{Recht van inzage}
    \begin{enumerate}
    \item Het Auditcomit\'e heeft het recht van inzage in elk document dat het nodig heeft om zijn opdracht uit te voeren. Zij behandelen deze documenten echter als strikt vertrouwelijk.
    \end{enumerate}

    \subsection{Rapportering}
    \begin{enumerate}
    \item Het Auditcomit\'e zal steeds een verslag van zijn vergaderingen ter beschikking stellen van de Algemene Vergadering.
    \end{enumerate}


\part{Werkgroepen}
\section{Definitie}
\begin{enumerate}
    \item Een werkgroep is een georganiseerde groep studenten aan de Faculteit Ingenieurswetenschappen die zich inzet voor de studenten aan de Faculteit Ingenieurswetenschappen.
    \item Een werkgroep wordt erkend door de Algemene Vergadering van VTK.
    \item Werkgroepen worden onderverdeeld in richtingsgebonden werkgroepen en niet-richtingsgebonden werkgroepen.
    \item De erkende richtingsgebonden werkgroepen zijn:
    \begin{itemize}
        \item \textcolor{green}{Biomedix}
        \item Chemix
        \item Existenz
        \item Mechanix
        \item Statix
    \end{itemize}
    \item De erkende niet-richtingsgebonden werkgroepen zijn:
    \begin{itemize}
        \item BEST
        \item \textcolor{red}{IAESTE}
        \item Revue
    \end{itemize}
    \item Een werkgroep is een onderdeel van VTK, maar beslist autonoom over hun interne werking. Onder interne werking vallen de dagelijkse operaties van de werkgroepen die nodig zijn voor het uitvoeren van de core activiteiten.
    \item Een richtingsgebonden werkgroep tracht een gevoel van eenheid te creëren binnen de richting op verschillende manieren, met een focus op het organiseren van verbredende evenementen zoals een bedrijfsbezoek, workshop, lezing of andere activiteiten in het vakgebied. Daarnaast zijn ontspannende activiteiten toegelaten indien ze duidelijk focussen op het samenbrengen van de specifieke doelgroep waar VTK dat niet kan. Deze activiteiten mogen niet de focus zijn van de wekrgroep en VTK kan met dit als reden ontspannende evenementen weigeren.

    Een niet-richtingsgebonden werkgroep focust op het verbreden van het scala aan activiteiten, ontwikkelingsmogelijkheden of ludieke evenementen die VTK zelf niet voorziet.
    \item Het oprichten van een werkgroep gebeurt door het voorleggen van een beleidsplan aan G5, dat erna wordt voorgelegd op een bestuursorgaan van VTK. Reeds bestaande werkgroepen en G5 geven advies en indien het bestuursorgaan goedkeuring geeft, zal de werkgroep op een AV worden voorgelegd ter goedkeuring.

    \item De werking van de werkgroepen worden elke 3 jaar geëvalueerd, namelijk elk derde werkingsjaar sinds de oprichting van de werkgroepen. Aa een evaluatiegroep, waarvan de leden worden goedgekeurd door het bestuursorgaan, wordt zijn taak gedelegeerd ter formulering van een advies naar de AV toe. Hier wordt ook op voorhand gecommuniceerd met de werkgroep in kwestie. De stichtingsjaren van de richtingsgebonden werkgroepenzijn als volgt:
    \begin{itemize}
        \item Existenz: 1995 (gedurende academiejaar 1994-1995)
        \item Statix: 2016 - 2017
        \item Mechanix: 2021 - 2022
        \item Chemix: 2022 - 2023
        \item Biomedix: 2023 - 2024
    \end{itemize}

    \item Dagelijksbestuurders van VTK kunnen niet aangesteld worden tot werkend lid ter vertegenwoordiging van een werkgroep.

\end{enumerate}

\section{Plichten}
\begin{enumerate}
    \item Per werkgroep zetelen er twee leden in de Algemene Vergadering. Ze worden dan ook aangesteld als werkend lid aan het begin van het werkingsjaar. Deze leden kunnen steeds van het Bestuursorgaan specifieke volmachten krijgen zoals bepaald in dit Huishoudelijk Reglement. 
    \item Werkgroepen moeten een beleidstekst ter bekrachtiging voorleggen aan de Algemene Vergadering op de eerste Algemene Vergadering na de overdracht. Hierin staat een korte beschrijving van de werkgroep, het te volgen beleid voor het komende werkingsjaar en een beschrijving van de manier waarop de opvolging op het einde van het jaar geregeld wordt. 
    \item Werkgroepen kunnen te allen tijde een herziening van deze beleidstekst aan het Bestuursorgaan of de Algemene Vergadering voorleggen. 
    \item  Werkgroepen moeten hun volledige begroting ter goedkeuring voorleggen op dezelfde Algemene Vergadering als VTK zijn begroting voorlegt ter goedkeuring. 
    \item Alle uitgaven van een werkgroep moeten goedgekeurd worden door het Bestuursorgaan van VTK.
    \item Wanneer een budget van een werkgroep door het Bestuursorgaan behandeld wordt, wordt een lid van de werkgroep hiervoor uitgenodigd.
    \item Activiteiten van werkgroepen worden vastgelegd in samenspraak met en onder goedkeuring van VTK.
    \item Contracten die nodig zijn voor activiteiten goedgekeurd op de begroeting, moeten door de Vice-Praeses in naam van de werkgroep getekend worden. Indien de Vice-Praeses niet tijdig reageert, mag dit ook door een ander bestuurslid.
    \item De plichten opgesomd in deze sectie zijn niet exhaustief. Het Bestuursorgaan, in samenspraak met de werkgroepen, kan altijd extra verplichtingen opleggen.
\end{enumerate}

\section{Rechten}
\begin{enumerate}
    \item Werkgroepen krijgen een rekeningnummer en bijhorende bankkaarten van VTK, beheerd door VTK, voor het begin van het werkingsjaar. 
    \item Werkgroepen mogen voor activiteiten die door het Bestuursorgaan goedgekeurd zijn, beroep doen op de verzekeringen van VTK. Eventuele extra benodigde verzekeringen worden door VTK afgesloten op vraag van de werkgroep. 
    \item Werkgroepen krijgen een vast bedrag teambuildingsbudget per persoon van VTK, dat bepaald wordt bij de begroting en in verhouding ligt met het budget van VTK praesidiumleden.
\end{enumerate}

\section{Subsidies}
\begin{enumerate}
    \item Werkgroepen kunnen een werkingssubsidie aanvragen bij VTK. Het is aan het Bestuursorgaan om het subsidiebudget van werkgroepen te bepalen.
    \item Een werkingssubsidie geldt voor het hele jaar en wordt niet gegeven per project of evenement.
    \item Indien werkingssubsidies worden aangevraagd, wordt er verwacht dat de werkgroep de richtlijnen van het bestuursorgaan betreffende de sponsoring per persoon en teambuildings volgt. Verder wordt verwacht dat de werkgroep VTK helpt door de dagelijkse werking te ondersteunen. Dit kan door 't Elixir en/ of Theokot open te houden. Dit doen ze tweemaal in het werkingsjaar. Aan het begin van het jaar kan de werkgroep een voorstel doen om een andere post te helpen bij het organiseren van een evenement in plaats hiervan.
    \item Ook al haalt de werkgroep meer geld binnen dan aanvankelijk gedacht, behoudt ze haar recht de werkingssubsidie te behouden onder voorwaarde dat deze ook effectief wordt besteed aan haar studenten. Dit wordt bepaald door het bestuursorgaan.
    \item Bij aanvang van een lustrumjaar kan een werkgroep per uitzondering een extra subsidie aanvragen bij VTK. Dit bedrag wordt bepaald door het Bestuursorgaan op de begrotingsvergadering voorafgaand aan het academiejaar.
    \item Indien een werkgroep een innovatief idee heeft, kan ze een projectsubsidie aanvragen. Deze wordt op dezelfde manier bepaald als de werkingssubsidie.
\end{enumerate}

\section{Ondersteuning}
\begin{enumerate}
    \item Werkgroepen mogen gebruikmaken van de logistieke diensten van VTK. 
    \item Er staat geen limiet op het gebruik van de voertuigen. De groepsco\"ordinator Logistiek en verantwoordelijke transport hebben echter het recht om onnodig gebruik te weigeren.
    \item Werkgroepen mogen geen eigen voorraad beheren. Er is één centrale plaats, het Flesserke genaamd, die beheerd wordt door VTK Logistiek en werkgroepen zijn verplicht om hiervan gebruik te maken. Overschotten moeten na een evenement dan ook steeds teruggebracht worden.
    \item Werkgroepen kunnen gebruikmaken van de diensten op Blok 6. Om onnodig gebruik van de printer te vermijden worden de printkosten opgenomen in de begroting. 
    \item Werkgroepen krijgen ruimte op de servers van VTK voor opslag en/of een website. Reeds bestaande domeinnamen kunnen alleen worden verwijderd of aangepast in samenspraak met de werkgroepen.
\end{enumerate}

\section{Samenwerking met bedrijven}
\begin{enumerate}
    \item  Richtingsgebonden werkgroepen mogen enkel contact opnemen met externen na expliciete toelating van de groepscoördinator van VTK Bedrijvenrelaties. Sponsorbrochures en stage brochures worden ook voorgelegd ter goedkeuring alvorens ze worden gebruikt. De geldigheid van de beslissing bij op een BO of AV besproken worden. Deze eenmalige toelating is geldig gedurende het lopende werkingsjaar. Hiernaast wordt VTK Bedrijvenrelaties steeds in cc gezet bij contact met bedrijven.
    \item Evenementen en publiciteit via elk mogelijk kanaal waar bedrijven het oog hebben op het rekruteren van studenten is ten allen tijde verboden. Indien er onduidelijkheid bestaat over de rekruteringsgerichte aard van een evenement of publicatie, kan dit ter beoordeling worden voorgelegd aan de Vice-Praeses en de Groepscoördinator BR voor goedkeuring.
\end{enumerate}

\section{Rekeningen}
\begin{enumerate}
    \item Een richtingsgebonden werkgroep krijgt elk jaar een bedrag toegewezen op hun rekening dat als cashflow fungeert en voldoende is om hun werking mogelijk te maken. Op het einde van het werkingsjaar dient dit bedrag terug op de rekening te staan.
    \item Voor de richtingsgebonden werkgroepen kan er geen geld overgedragen worden van jaar op jaar. Eventuele verliezen worden opgevangen door VTK indien er hiervoor een adequate verantwoording wordt gegeven door de richtingsgebonden werkgroep. Dit houdt in dat de voorzitter en beheerder van een werkgroep op de laatste BO/AV van het werkingsjaar uitleg en verantwoording dienen te geven bij het resultaat van hun werkingsjaar. Dit laatste geldt ook voor niet-richtingsgebonden werkgroepen.
\end{enumerate}

\section{Communicatie}
\begin{enumerate}
    \item VTK fungeert als het centrale aanspreekpunt voor studenten aan de faculteit, vertegenwoordigd door de Praeses en de Groepscoördinator Onderwijs. Overleg tussen departementen en richtingsgebonden werkgroepen zonder dat de Praeses, Vice-Praeses en Groepscoördinator Onderwijs op de hoogte zijn, is niet toegestaan. Indien een werkgroep communicatie met een professor of het departement wenst op te starten over een specifiek onderwerp, dient vooraf toestemming te worden gevraagd aan de Vice-Praeses en de Groepscoördinator Onderwijs. Na goedkeuring kan de communicatie direct verlopen, mits de Vice-Praeses en de Groepscoördinator Onderwijs steeds in CC worden opgenomen. Deze eenmalige toestemming is geldig voor het lopende academiejaar.
    \item Er wordt geen werkgroepgerelateerde materie besproken op departementale vergaderingen. De leden van de departementale vergadering zitten daar enkel en alleen als richtingsvertegenwoordiger.
    \item Richtingsgebonden werkgroepen mogen evenementen organiseren samen met het departement, indien de Vice-Praeses van alle plannen op de hoogte is en hiermee akkoord gaat.
    \item De Vice-Praeses krijgt op 15 juli toegang als beheerder op de sociale media van de richtingsgebonden werkgroepen.
    \item In alle communicatie van de richtingsgebonden werkgroepen naar externen wordt duidelijk geschreven dat de werkgroep een onderdeel is van VTK. Ook op de website van de werkgroep wordt een vermelding van VTK gemaakt.
    \item In het logo van de richtingsgebonden werkgroep komt een duidelijke vermelding van VTK te staan. Dit dient jaarlijks goedgekeurd te worden door een bestuurslid.
    \item De werkgroepen hebben jaarlijks het recht een regeling te treffen met VTK Communicatie om gebruik te maken van de communicatiekanalen van VTK. Dit gebeurt aan het begin van het werkingsjaar. De afzonderlijke regelingen worden transparant en met motivatie naar alle werkgroepen gecommuniceerd.
\end{enumerate}

\section{Verkiezingen en opvolging}
\begin{enumerate}
    \item Een richtingsgebonden werkgroep is verplicht open verkiezingen te organiseren, waarbij alle studenten van de betrokken opleiding(en) welkom zijn om deze bij te wonen of zich kandidaat te stellen. De verkiezing is open voor waarnemers van de hele faculteit.
    \item Niet-richtingsgebonden werkgroepen zijn niet verplicht verkiezingen te organiseren, maar worden geacht zo inclusief mogelijk te zijn bij de opstart van hun werking.
\end{enumerate}

\part{Draaiboek van de vergaderingen}
\section{Algemeen Draaiboek}\label{sec:AlgemeenDraaiboek}

    \begin{enumerate}
    \item De vergadervoorzitter blijft neutraal gedurende de stemming tenzij bij onbesliste stemming. Indien een stemming onbeslist is, is de stem van de vergadervoorzitter beslissend.
    \item De vergadervoorzitter mag geen volmacht dragen.
    \item Beslissingen worden genomen bij consensus. Indien er geen consensus bereikt wordt, dient er gestemd te worden. Elk lid van de vergadering kan de stemming aanvragen.
    \item Iemand kan niet aan een stemming over een agendapunt deelnemen, indien hij/zij niet bij de volledige discussie over dat agendapunt aanwezig was.
    \item Een stemming gebeurt in principe bij handopsteking, behalve indien anders vermeld in dit reglement. Geheime stemming kan echter ook door \'e\'en van de aanwezige leden aangevraagd worden. Een betrokken natuurlijke persoon woont in geen geval de stemming bij.
    \item Op iedere vergadering van VTK gebeurt de stemming volgens de principes in dit reglement, tenzij anders vermeld.
    \item Indien er \'e\'en voorstel voorligt, wordt er gestemd bij absolute meerderheid. Indien er geen absolute meerderheid voor of tegen is, wordt de discussie heropend als minstens \'e\'en aanwezig lid daarom vraagt, en wordt er vervolgens opnieuw gestemd waarbij de onthoudingen niet meetellen.
    \item Een stemming over \'e\'en voorstel is onbeslist als er bij de tweede stemming over hetzelfde voorstel evenveel voor als tegenstemmen zijn (stemmen voor = stemmen tegen).
    \item Indien er meerdere voorstellen voorliggen (zowel ordinair als niet-ordinair), beslist de vergadervoorzitter over de te volgen stemprocedure.
    \end{enumerate}

% To Do: Naar HR
\section{De Grocovergadering}

    \begin{enumerate}
    \item De Grocovergadering wordt bijeengeroepen door de Praeses. Tevens is de Praeses verplicht de Grocovergadering samen te roepen wanneer minstens \'e\'en vijfde van de leden van de Grocovergadering hierom schriftelijk of elektronisch verzoekt. De datum waarop deze Grocovergadering zal plaatsvinden en de agenda wordt dan in overleg met de betrokken leden door de Praeses bepaald.
    \item Voor stemmingen op de Grocovergadering geldt de neutraliteit van de Praeses niet. 
    \item De leden van de Grocovergadering kunnen zich niet laten vertegenwoordigen door middel van volmachten.
    \item Personen die niet van ambtswege op de Grocovergadering aanwezig moeten zijn, kunnen, in overleg met de vergadervoorzitter, \'e\'en of verscheidene agendapunten van de vergadering bijwonen.
    \item De notulen worden genomen door een lid van de Grocovergadering en worden ter beschikking gesteld van het Bestuursorgaan.
    \item Als een Groco zelf niet aanwezig kan zijn op de Grocovergadering kan hij/zij één  iemand van zijn team laten deelnemen in zijn plaats in overleg met de vergadervoorzitter en na instemming van de Grocovergadering.
    \end{enumerate}

\section{Het Bestuursorgaan}
    \begin{enumerate}
    \item Het Bestuursorgaan komt in principe elke week op zondag bijeen en telkens wanneer het doel of het belang van VTK het eist. Tevens is de Praeses verplicht het Bestuursorgaan samen te roepen wanneer minstens \'e\'en vijfde van de leden van het Bestuursorgaan hierom schriftelijk of elektronisch verzoekt. De datum waarop deze vergadering zal plaatsvinden en de agenda wordt dan in overleg met de betrokken leden door de Grocovergadering bepaald.
    \item De leden van het Bestuursorgaan kunnen zich niet laten vertegenwoordigen door middel van volmachten.
    \item Personen die niet van ambtswege op het Bestuursorgaan aanwezig moeten zijn, kunnen, in overleg met de vergadervoorzitter en na instemming van het Bestuursorgaan, \'e\'en of verscheidene agendapunten van de vergadering bijwonen.
    \item Het Bestuursorgaan kan beslissen dat een vergadering of een deel van een vergadering waarop deze personen aanwezig zijn, gesloten wordt verklaard, d.w.z. enkel bijgewoond kan worden door de leden van de vergadering.
    \end{enumerate}

\part{Studentenvertegenwoordiging}
    \section{Mandaten}
    \begin{enumerate}
    \item Bij uitoefening van mandaten als studentenvertegenwoordiger aan de faculteit, universiteit, associatie of daarbuiten is men geen verantwoording verschuldigd aan geen enkel orgaan of persoon binnen VTK indien het vertrouwelijke dossiers betreft, tot het moment dat deze vertrouwelijkheid niet meer vereist is. 
    \end{enumerate}

\section{Het VTK Bureau}
\textbf{Samenstelling}
    \begin{enumerate}
    \item In het VTK Bureau zetelen alle studenten van de Faculteit Ingenieurswetenschappen, zoals gedefinieerd in de statuten van VTK. 
    \item De Groco Onderwijs zit het VTK Bureau voor of stelt bij afwezigheid een vervanger aan. Indien er geen vervanger is aangesteld, wordt de vergadering opgeschort.
    \item Andere personen kunnen, in overleg met de vergadervoorzitter, één of verscheidene agendapunten van de vergadering bijwonen, maar mogen niet deelnemen aan de stemming.
    \end{enumerate}

    \textbf{Doel en bevoegdheden}
    \begin{enumerate}
    \item Het VTK Bureau geldt als de open onderwijsvergadering van VTK. Op deze vergaderingen kunnen onderwerpen uit alle vergaderingen van de departementen, de faculteit, de groep, de campus en de universiteit besproken worden. 
    \item Het VTK Bureau geeft advies aangaande alle onderwerpen vermeld onder paragraaf 1 van dit artikel.
    \item Het verslag van het VTK Bureau wordt na goedkeuring op het VTK Bureau, in principe op de eerstkomende vergadering van het Bestuursorgaan voorgelegd ter informatie.
    \item Alle beslissingen van het VTK Bureau worden aan het Bestuursorgaan ter informatie voorgelegd en besproken onder leiding van de Groco Onderwijs. 
    \item Over standpunten die niet op het VTK Bureau besproken zijn, beslissen de Praeses en de leden van de post Onderwijs, indien mogelijk aangevuld met de leden van de VTK-delegatie op de desbetreffende vergadering. Deze beslissing wordt toegelicht op het eerstvolgende VTK Bureau.
    \end{enumerate}

    \textbf{Besluitvorming}\label{art:bureau-besluit}
    \begin{enumerate}
    \item In het geval strikte vertrouwelijkheid vereist is, mag het FO-bestuur, bestaande uit de leden van de post Onderwijs en de Praeses, een beslissing nemen. Zij zijn hierover geen verantwoording verschuldigd, noch naar het Bestuursorgaan, noch naar het VTK Bureau tot op het moment dat de vertrouwelijkheid niet meer vereist is.
    \end{enumerate}

    \textbf{Openbaarheid}
    \begin{enumerate}
    \item De voorzitter van het VTK Bureau kan beslissen dat een vergadering of een deel ervan gesloten wordt verklaard, d.w.z. enkel bijgewoond kan worden door de student-POC'ers, Praeses en de post Onderwijs.
    \end{enumerate}

% To Do: naar HR, aanpassen na participatiereglement
\section{RiVer-verkiezingen}\label{sec:vertegenwoordiging-verkiezing}

    \textbf{Doel}
    \begin{enumerate}
    \item De hieronder beschreven procedure dient ter verkiezing van de studentendelegatie in de verschillende opleidings-POC's, departementsraden en -besturen. 
    \item Zolang de procedure niet is afgerond blijft de studentendelegatie van het vorige academiejaar actief in zoverre zij voldoen aan voorwaarden bepaald in de statuten van VTK.
    \end{enumerate}
    
    \textbf{Kandidaatstelling}
    \begin{enumerate}
    \item Aan het begin van het academiejaar doet VTK een oproep waarbij studenten aan de Faculteit Ingenieurswetenschappen zich kandidaat kunnen stellen voor (een) positie(s) in de verschillende opleidings-POC's, departementsraden en -besturen tot de deadline voor kandidaatstelling.
    \item Na de deadline voor kandidaatstelling worden de namen van de verschillende kandidaten bekendgemaakt op de website van VTK.
        \begin{itemize}
        \item Indien er te weinig of net genoeg kandidaten zijn, wordt de procedure zoals in artikel \ref{art:vertegenwoordiging-stilzwijgend} gevolgd.
        \item Indien er te veel kandidaten zijn, wordt er een compromis gezocht van waarnemers en stemgerechtigden, anders wordt de procedure zoals in artikel \ref{art:vertegenwoordiging-schriftelijk} gevolgd.
        \end{itemize}
    \item De deadline voor kandidaatstelling wordt bepaald door de post Onderwijs.
    \end{enumerate}
    
    \textbf{Stilzwijgende procedure}\label{art:vertegenwoordiging-stilzwijgend}
    \begin{enumerate}
    \item Studenten kunnen bezwaar aantekenen tegen een kandidaat tot de deadline voor bezwaar aantekenen.
    \item Na de deadline voor bezwaar aantekenen worden
        \begin{itemize}
        \item de kandidaten zonder aangetekend bezwaar verkozen als studentenvertegenwoordiger voor het desbetreffende orgaan.
        \item de kandidaten met aangetekend bezwaar voorgelegd volgens de procedure in artikel \ref{art:vertegenwoordiging-schriftelijk}.
        \end{itemize}
    \item De deadline voor bezwaar aantekenen wordt bepaald door de post Onderwijs.
    \end{enumerate}
    
    \textbf{Schriftelijke procedure}\label{art:vertegenwoordiging-schriftelijk}
    \begin{enumerate}
    \item VTK organiseert een schriftelijke stemming voor de kandidaten waarbij volgende studenten stemgerechtigd zijn:
        \begin{itemize}
        \item alle studenten uit de desbetreffende opleiding in het geval van een POC.
        \item alle studenten uit opleidingen die onder het desbetreffende departement vallen in het geval van een departementsraad of -bestuur. 
        \end{itemize}
    \item Kandidaten worden met een gewone meerderheid na afloop van de stemming verkozen.
    \end{enumerate}

\section{Samenstelling VTK-studentenfractie}

    \textbf{Facultair}
    \begin{enumerate}
    \item Voor de samenstelling van de studentenfractie in de faculteitsraad en in de FPOC legt de Groco Onderwijs een lijst ter bekrachtiging voor aan het Bestuursorgaan.
    \item De samenstelling van de studentendelegaties in de opleidings-POC’s, departementsraden en departementsbesturen wordt beschreven in Afdeling \ref{sec:vertegenwoordiging-verkiezing}.
    \item De Praeses en de Groco Onderwijs vertegenwoordigen de studenten in het faculteitsbestuur.
    \end{enumerate}

    \textbf{Campus}
    \begin{enumerate}
    \item  De samenstelling van de stemgerechtigde VTK-delegatie op de LOKO AV wordt beslist door de Praeses.
    \item De samenstelling van de VTK-delegatie in het geval van een gesloten LOKO AV wordt beslist door de Praeses.
    \item Bij afwezigheid van de Praeses, worden de bevoegdheden in §1 en §2 uitgeoefend door de Groco Onderwijs. Indien beiden afwezig zijn, wordt een vervanger aangeduid.
    \end{enumerate}

    \textbf{Universitair}
    \begin{enumerate}
    \item  De samenstelling van de stemgerechtigde VTK-delegatie op de Stura AV wordt beslist door de Groco Onderwijs%, rekening houdend met de volgende volgorde van voorrang: De Praeses en de Groco Onderwijs, de leden van de post Onderwijs, de aanwezige praesidiumleden, aanwezige POC’ers, studenten van de faculteit.
    \item De samenstelling van de VTK-delegatie in het geval van een gesloten Stura AV wordt beslist door de Groco Onderwijs.
    \item Bij afwezigheid van de Groco Onderwijs, worden de bevoegdheden in §1 en §2 uitgeoefend door de Praeses. Indien beiden afwezig zijn, wordt een vervanger aangeduid.
    \end{enumerate}

\part{Diverse bepalingen}
\section{Volmachten}
    \begin{enumerate}
    \item De penningmeester(s) krijgen een volmacht op de rekeningen van de vereniging. De penningmeester(s) hebben een volmacht om, conform dit Huishoudelijk Reglement, volmachten voor specifieke rekeningen van de vereniging uit te delen aan andere bestuurders en de werkgroepen.
    \item De voorzitter van de vereniging krijgt een volmacht op de rekening(en) van de vereniging. 
    \end{enumerate}

\section{Historiek der wijzigingen}
\begin{enumerate}
    \item Dit Huishoudelijk Reglement werd aangenomen op de AV van de Vlaamse Technische Kring vzw op 01 juli 2024.
    \item Dit huishoudelijk regelement werd gewijzigd op de AV van Vlaamse Technische Kring vzw op 30 juni 2025.
    \item Dit huishoudelijk regelement werd gewijzigd op de AV van Vlaamse Technische Kring vzw op 6 juli 2026.
\end{enumerate}

% \part{Praesidiumverkiezingen}\label{TitelVerkiezingen}

%     \section{Verkiezingen}\label{ArtikelVerkiezingen}
%     \begin{enumerate}
%     \item Het Bestuursorgaan is verplicht elk jaar tussen februari en juni verkiezingen te organiseren, de ``praesidiumverkiezingen''\ genaamd. Op deze verkiezingen zijn alle studenten aan de Faculteit Ingenieurswetenschappen aan de KU Leuven stemgerechtigd.
%     \item De precieze modaliteiten van de verkiezingen moeten worden vastgelegd in een specifiek verkiezingsreglement. Het Bestuursorgaan keurt het verkiezingsreglement goed. Het Bestuursorgaan moet het verkiezingsreglement openbaar maken op de website van de vereniging. Een aankondiging van deze publicatie dient te verschijnen in een papieren publicatie die gratis ter beschikking van de studenten aan de Faculteit Ingenieurswetenschappen gesteld wordt. Vanaf de bekendmaking van VKK (verkiezingscontrolecommissie) t.e.m. de dag van de jaarlijkse verkiezing mag het reglement geen wijzigingen ondergaan.
%     \item De organisatie van de verkiezingen verloopt eveneens met het oog op de verkiezingen van de vertegenwoordigers van de studenten bij de KU Leuven. Daarbij respecteert het Bestuursorgaan het Participatiedecreet en het Participatiereglement van Studentenraad KU Leuven en LOKO.
%     \item Deelnemen aan de verkiezingen dient in ploegverband te gebeuren. Een deelnemende ploeg moet mimimum vijf leden tellen. De leden van de verkiesbare ploeg moeten studenten aan de Faculteit Ingenieurswetenschappen zijn. Ten laatste twee weken voor de verkiezingsdag moet het Bestuursorgaan op de website van de vereniging de lijst met deelnemende ploegen publiceren, met inbegrip van de namen van de leden van alle ploegen.
%     \item Het verkiezingsreglement dient in elk geval te specificeren op welke manier bepaald wordt welke ploeg de verkiezingen heeft gewonnen. Er kan slechts één ploeg de verkiezingen winnen.
%     \end{enumerate}
    
%     \section{Bekrachtiging van de uitslag van de verkiezingen}
%     \begin{enumerate}
%     \item Het Bestuursorgaan ziet er op toe dat de bepalingen van artikel \ref{ArtikelVerkiezingen} van deze statuten en van het volledige verkiezingsreglement worden nageleefd. Vervolgens moet het Bestuursorgaan, na controle van het bepaalde in artikel \ref{ArtikelVerkiezingen} en in het verkiezingsreglement, beslissen of de verkiezingen conform de statuten en het verkiezingsreglement hebben plaatsgevonden. De bekrachtiging van de verkiezingsuitslag houdt in dat het Bestuursorgaan aanvaardt dat de leden van de winnende ploeg op de laatste Algemene Vergadering van het werkingsjaar toetreden als werkend lid.
%     \item Indien zij deze verplichting verzuimt, komt het de dagelijks bestuurders die het komende academiejaar niet langer student aan de Faculteit Ingenieurswetenschappen zullen zijn toe een aantal personen aan te wijzen die als werkend lid tot de vereniging zullen worden toegelaten, rekening houdend met het bepaalde in artikel \ref{ArtikelWerkendeLeden}. Deze dagelijks bestuurders kunnen discretionair en zonder motivatie de begin- en einddatum van het lidmaatschap van die personen bepalen er rekening mee houdende dat het lidmaatschap maximaal 1 jaar kan duren.  Deze personen staan in voor de continuïteit van de werking van VTK en zijn verplicht om nieuwe verkiezingen te organiseren.
%     \end{enumerate}

% \part{Het Bestuursorgaan}

%     \section{Bestuurders}
%     \begin{enumerate}
%     \item Het Bestuursorgaan kiest onder de bestuurders een voorzitter (indachtig de voorwaarden in \ref{ParVoorwaardeVoorzitter} van dit artikel), ondervoorzitter, secretaris en minstens één penningmeester.
%     \item Het Bestuursorgaan zal bij het kiezen van een voorzitter rekening houden met volgende voorwaarde:
%         \begin{itemize}
%         \item De voorzitter moet in het jaar dat zijn mandaat van kracht is, zijn bachelordiploma aan de Faculteit Ingenieurswetenschappen behaald hebben en ingeschreven zijn voor een masteropleiding aan de Faculteit Ingenieurswetenschappen.
%         \end{itemize}\label{ParVoorwaardeVoorzitter}
%     \item De bestuurders voeren hun mandaat onbezoldigd uit.
%     \end{enumerate}
    
%     \section{Einde van het bestuurderschap}
%     \begin{enumerate}
%     \item Een bestuurder is verplicht om na zijn ontslag zijn opdracht verder te vervullen totdat redelijkerwijs in zijn vervanging kan worden voorzien voor een maximum van 1 maand. 
%     \end{enumerate}
    
%     \section{Bijeenkomst van het Bestuursorgaan}
%     \begin{enumerate}
%     \item Het Bestuursorgaan vergadert na oproeping door de voorzitter zo dikwijls als het belang van de vereniging het vereist. De vergadering wordt gehouden op de plaats die in het oproepingsbericht wordt vermeld.
%     \item De oproeping dient minstens 24 uur voor aanvang van de vergadering te gebeuren.
%     \item Het Bestuursorgaan wordt voorgezeten door de voorzitter. Indien deze niet aanwezig kan zijn stelt hij een vervanger aan. Deze vervanger wordt vermeld in de uitnodiging. In het geval er geen vervanger aangesteld is wordt het Bestuursorgaan voorgezeten door de ondervoorzitter. Indien de ondervoorzitter eveneens afwezig is, wordt het Bestuursorgaan voorgezeten door de oudst aanwezige dagelijks bestuurder. Indien deze allen afwezig zijn, beslist het Bestuursorgaan onder leiding van de oudst aanwezige bestuurder of de bijeenkomst wordt uitgesteld.
%     \item Het Bestuursorgaan kan slechts geldig vergaderen wanneer ten minste 1/4 van zijn leden aanwezig zijn. Het Bestuursorgaan oefent zijn bevoegdheden uit als college.
%     \item Van elke vergadering worden notulen opgesteld, ondertekend door de voorzitter en de secretaris van het Bestuursorgaan. De notulen kunnen tevens getekend worden door de bestuurders die het wensen. Ze worden bewaard in een notulenregister dat ter inzage zal zijn van de werkende leden, die hun inzagerecht zullen uitoefenen overeenkomstig de WVV-wet vastgelegde modaliteiten.
%     \end{enumerate}
    
%     \section{Belangenvermenging}
%     \begin{enumerate}
%     \item Indien een bestuurder, rechtstreeks of onrechtstreeks, voor een beslissing die onder de bevoegdheid valt van het Bestuursorgaan een belang van vermogensrechtelijke of persoonlijke aard heeft dat strijdig is met het belang van de vereniging, moet hij dit meedelen aan de andere bestuurders voor het Bestuursorgaan een besluit neemt. De verklaring van deze bestuurder over de aard van het strijdig belang wordt opgenomen in de notulen van de vergadering.
%     \item Het Bestuursorgaan kan met gewone meerderheid van de stemmen beslissen of iemand al dan niet een belangenconflict van vermogensrechtelijke of persoonlijke aard heeft.
%     \item De bestuurder met het tegenstrijdig belang verwijdert zich uit de vergadering en onthoudt zich van de beraadslaging en de stemming over de aangelegenheid waarop het betrekking heeft.
%     \item Voormelde procedure is niet toepasselijk op gebruikelijke verrichtingen die plaatshebben onder de voorwaarden en tegen de zekerheden die op de markt gewoonlijk gelden voor soortgelijke verrichtingen.
%     \item Indien een bestuurder lid is van een (niet-)verkiesbare ploeg verwijdert hij zich uit de vergadering indien er een beraadslaging plaatsvindt die betrekking heeft op de controle op de verkiezingen of de bespreking van het verkiezingsresultaat.
%     \end{enumerate}
    
%     \section{Bevoegdheden van het Bestuursorgaan}
%     \begin{enumerate}
%     \item Het Bestuursorgaan is bevoegd om alle handelingen van intern bestuur te verrichten die nodig of dienstig zijn tot verwezenlijking van het doel van de vereniging, met uitzondering van die handelingen waarvoor volgens de wet of deze statuten de Algemene Vergadering exclusief bevoegd is.
%     \item Het Bestuursorgaan beslist over het Intern Reglement bij gewone meerderheid van de stemmen van de aanwezige of vertegenwoordigde werkende leden.
%     \item Onverminderd de verplichtingen die voortvloeien uit collegiaal bestuur, met name overleg en toezicht, kunnen de bestuurders de bestuurstaken onder elkaar verdelen. Zodanige verdeling van taken kan aan derden niet worden tegengeworpen, zelfs niet nadat ze zijn openbaar gemaakt. Niet-naleving ervan brengt wel de interne aansprakelijkheid van de betrokken bestuurder(s) in het gedrang.
%     \item Het Bestuursorgaan kan een deel van zijn bestuursbevoegdheid delegeren aan één of meerdere derden niet-bestuurders, zonder dat deze overdracht evenwel betrekking kan hebben op het algemeen beleid van de vereniging of de algemene bestuursbevoegdheid van het Bestuursorgaan. Hierin wordt een onderscheid gemaakt tussen vertegenwoordigers van werkgroepen en bijzondere mandaten. De bevoegdheid van vertegenwoordigers van werkgroepen wordt beschreven in het Intern Reglement, die van andere mandaten wordt precies afgebakend door het Bestuursorgaan, die ook de duur van het mandaat bepaalt. Het mandaat kan te allen tijde met onmiddellijke ingang door het Bestuursorgaan worden ingetrokken.
%     \end{enumerate}

%     \section{Vertegenwoordiging}\label{ArtikelVertegenwoordiging}
%     \begin{enumerate}
%     \item Het Bestuursorgaan vertegenwoordigt als college de vereniging in alle handelingen in en buiten rechte.
%     \item Onverminderd de algemene vertegenwoordigheidsbevoegdheid van het Bestuursorgaan als college, wordt de vereniging in en buiten rechte eveneens geldig vertegenwoordigd door ofwel:
%         \begin{itemize}
%         \item twee dagelijks bestuurders, waarvan één van beide voorzitter, ondervoorzitter, penningmeester of secretaris is; ofwel
%         \item één dagelijks bestuurder en één bestuurder tot een bedrag van maximaal EUR 7500,-; ofwel
%         \item één dagelijks bestuurder tot een bedrag van EUR 2500,- voor wat het dagelijks bestuur betreft dat binnen de bevoegdheid valt van deze persoon zoals beschreven in het Intern Reglement.
%         \end{itemize}\label{ParLimietenVertegenwoordiging}
%     \item Het Bestuursorgaan kan, wanneer de noodzakelijkheid en het belang van de vereniging zulks vereisen, personen uitzonderlijk machtiging geven om ten aanzien van derden te handelen tot een door het Bestuursorgaan gespecificeerd bedrag.
%     \end{enumerate}
    
%     \section{Dagelijks bestuur}
%     \begin{enumerate}
%     \item Het dagelijks bestuur van de vereniging alsook de vertegenwoordiging wat betreft dat bestuur kan door het Bestuursorgaan worden opgedragen aan één of meer van haar leden.
%     \item De dagelijks bestuurder(s) is (zijn) wat betreft externe vertegenwoordigingsmacht bevoegd te handelen volgens de bepalingen in artikel \ref{ArtikelVertegenwoordiging}, \ref{ParLimietenVertegenwoordiging}.
%     \item Het dagelijks bestuur omvat zowel de handelingen en de beslissingen die niet verder reiken dan de behoeften van het dagelijks leven van de vereniging, als de handelingen en de beslissingen die, ofwel om reden van hun minder belang dat ze vertonen, ofwel omwille van hun spoedeisend karakter, de tussenkomst van het Bestuursorgaan niet rechtvaardigen.
%     \item De benoeming van de personen belast met het dagelijks bestuur en hun ambtsbeëindiging wordt openbaar gemaakt door neerlegging in het verenigingsdossier ter griffie van de ondernemingsrechtbank, en van een uittreksel daarvan bestemd om in de Bijlagen bij het Belgisch Staatsblad te worden bekendgemaakt. Uit de stukken moet in ieder geval blijken of de personen die vereniging vertegenwoordigen inzake dagelijks bestuur, de vereniging ieder afzonderlijk, gezamelijk, dan wel als college verbinden alsook de omvang van hun bevoegdheden.
%     \item In uitzonderlijke gevallen, wanneer de dringende noodzakelijkheid en het belang van de vereniging zulks vereisen, kunnen de besluiten van het Bestuursorgaan worden genomen bij eenparig schriftelijk akkoord van de dagelijks bestuurders. Daartoe is vereist dat er vooraf een unaniem akkoord is onder de dagelijks bestuurders om tot schriftelijke besluitvorming over te gaan. Schriftelijke besluitvorming veronderstelt in elk geval dat er een beraadslaging plaatsvond per e-mail, video- of telecommunicatie.
%     \end{enumerate}

% \part{Diverse bepalingen}

%     \section{Volmachten}
%     \begin{enumerate}
%     \item De penningmeester(s) krijgen een volmacht op de rekeningen van de vereniging. De penningmeester(s) hebben een volmacht om, conform het Intern Reglement, volmachten voor specifieke rekeningen van de vereniging uit te delen aan andere bestuurders en de werkgroepen.
%     \item De voorzitter van de vereniging krijgt een volmacht op de rekening(en) van de vereniging. 
%     \end{enumerate}
    
%     \section{Historiek der wijzigingen}
%     \begin{enumerate}
%     \item Deze geco\"ordineerde statuten werden aangenomen op de AV van 't ElixIr vzw op 3 november 2002.
%     \item Deze geco\"ordineerde statuten werden gewijzigd op de AV van 't ElixIr vzw op 1 juli 2005.
%     \item Deze geco\"ordineerde statuten werden gewijzigd op de AV van 't ElixIr vzw op 12 november 2006.
%     \item Deze geco\"ordineerde statuten werden gewijzigd op de AV van 't ElixIr vzw op 2 september 2007.
%     \item Deze geco\"ordineerde statuten werden gewijzigd op de AV van 't ElixIr vzw op 26 augustus 2008.
%     \item Deze geco\"ordineerde statuten werden gewijzigd op de AV van VTK Ondersteuning vzw op 16 november 2008.
%     \item Deze geco\"ordineerde statuten werden gewijzigd op de AV van VTK Ondersteuning vzw op 28 november 2010.
%     \item Deze geco\"ordineerde statuten werden gewijzigd op de AV van VTK Ondersteuning vzw op 22 oktober 2012.
%     \item Deze geco\"ordineerde statuten werden gewijzigd op de AV van VTK Ondersteuning vzw op 8 november 2015.
%     \item Deze geco\"ordineerde statuten werden gewijzigd op de AV van VTK Ondersteuning vzw op 1 mei 2016.
%     \item Deze geco\"ordineerde statuten werden gewijzigd op de AV van VTK Ondersteuning vzw op 19 mei 2021.
%     \item Deze geco\"ordineerde statuten werden gewijzigd op de AV van VTK Ondersteuning vzw op 1 juli 2021.
%     \item Deze geco\"ordineerde statuten werden gewijzigd op de AV van VTK Ondersteuning vzw op 14 juli 2022. 
%     \item Deze geco\"ordineerde statuten werden gewijzigd op de AV van Vlaamse Technische Kring vzw op 3 juli 2023.
%     \end{enumerate}

\end{document}